\documentclass[10pt]{article}
\usepackage[T1]{fontenc}
\usepackage[utf8]{inputenc}
\usepackage[french]{babel}
\usepackage{amsmath,amssymb,mathrsfs,mathpazo}
\usepackage{geometry}
\geometry{a4paper,margin=2cm}
\usepackage{graphicx,array,tabularx,tabularray,colortbl,multirow,multicol,booktabs,enumitem,pifont,ulem,fancybox,diagbox,cancel}
\usepackage{pstricks,pstricks-add,pst-plot,pst-tree,pst-node,pst-text,pst-func,pst-math,pst-bspline,pst-3dplot,pst-eucl,pst-solides3d}
\let\pscalculate\undefined
\usepackage{pst-bezier}
\usepackage{tikz,pgfplots,tkz-tab}
\usetikzlibrary{arrows,arrows.meta,patterns,calc,positioning,trees,decorations.markings,shapes.geometric,shapes.misc,shapes.arrows,intersections,decorations.pathmorphing,decorations.pathreplacing,backgrounds,shadows,babel}
\pgfplotsset{compat=1.18}
\usepgfplotslibrary{fillbetween}
\usepackage[np]{numprint}
\usepackage{listings,makecell,pas-tableur,twoopt}
\let\barre\undefined
\usepackage[pstricks]{tablvar}
\definecolor{darkgreen}{RGB}{0,100,0}
\providecommand{\red}{\color{red}}
\providecommand{\blue}{\color{blue}}
\providecommand{\green}{\color{green}}
\providecommand{\black}{\color{black}}
\providecommand{\white}{\color{white}}
\providecommand{\gray}{\color{gray}}
\providecommand{\cyan}{\color{cyan}}
\providecommand{\magenta}{\color{magenta}}
\providecommand{\yellow}{\color{yellow}}
\providecommand{\violet}{\color{violet}}
\providecommand{\orange}{\color{orange}}
\providecommand{\euro}{\texteuro}
\providecommand{\EUR}{\texteuro}
\providecommand{\ieme}{\textsuperscript{e}}
\providecommand{\degres}{\ensuremath{^\circ}}
\providecommand{\degree}{\ensuremath{^\circ}}
\providecommand{\vv}[1]{\overrightarrow{#1}}
\providecommand{\Coord}[1]{#1}
\providecommand{\decosix}{}
\providecommand{\starredbullet}{\ensuremath{\bullet}}
\pagestyle{empty}
\setlength{\parindent}{0pt}
\usepackage{adjustbox}\DeclareUnicodeCharacter{2236}{:}\DeclareUnicodeCharacter{2190}{\ensuremath{\leftarrow}}
\providecommand{\R}{\mathbb{R}}
\providecommand{\N}{\mathbb{N}}
\providecommand{\Z}{\mathbb{Z}}
\providecommand{\Q}{\mathbb{Q}}
\providecommand{\D}{\mathbb{D}}
\providecommand{\C}{\mathbb{C}}
\providecommand{\e}{{\mathrm{e}}}
\providecommand{\vect}[1]{\overrightarrow{#1}}
\providecommand{\vectt}[1]{\overrightarrow{\mathrm{#1}}}
\providecommand{\vvt}[1]{\overrightarrow{\mathrm{#1}}}
\providecommand{\np}[2][]{{#2~#1}}
\providecommand{\vv}[1]{{\overrightarrow{#1}}}
\renewcommand{\d}{{\mathrm{d}}}
\providecommand{\degres}{{^\circ}}
\providecommand{\degre}{{^\circ}}
\providecommand{\blue}{\color{blue}}
\providecommand{\red}{\color{red}}
\providecommand{\enskip}{\quad}
\providecommand{\llbracket}{\lbrack\!\lbrack}
\providecommand{\rrbracket}{\rbrack\!\rbrack}
\providecommand{\euro}{{\text{€}}}
\providecommand{\no}{{\text{n°}}}
\providecommand{\interval}[2]{{[#1;#2]}}
\providecommand{\intervFO}[2]{{[#1;#2[}}
\providecommand{\dotfill}{\ldots}

\providecommand{\ds}{\displaystyle}
\providecommand{\textding}[1]{#1}
\providecommand{\Oij}{\(\mathrm{O};\vec{\imath},\vec{\jmath}\)}
\providecommand{\Oijk}{\(\mathrm{O};\vec{\imath},\vec{\jmath},\vec{k}\)}
\providecommand{\Ouv}{\(\mathrm{O};\vec{u},\vec{v}\)}
\providecommand{\intervFF}[2]{\ensuremath{[#1;#2]}}
\renewcommand{\euro}{\ifmmode\text{\texteuro}\else\texteuro\fi}
\DeclareUnicodeCharacter{00B0}{\ifmmode^{\circ}\else\textdegree\fi}
\DeclareUnicodeCharacter{20AC}{\ifmmode\text{\texteuro}\else\texteuro\fi}
% Énoncés : Éducation nationale. Transcription et corrigés : APMEP, Denis Vergès et contributeurs. Découpage et classement : CoopMaths, Sébastien Lozano, Éric Elter et contributeurs. Adaptation de la présentation : Maths Explorer.
% CC BY-SA : https://coopmaths.fr/www/about/
% https://coopmaths.fr/alea/static/e3c/2020/tex/e3c_2020_00_sujet27_3_cor.tex
\begin{document}
\noindent\textit{APMEP / CoopMaths. CC BY-SA. Présentation adaptée par Maths Explorer.}\par\medskip
\subsection*{1.}
	
	\noindent
	\textbf{--- Premier modèle :}
	
	On fait l'hypothèse que ce nombre augmente de 21 \% par an. On définit alors une suite \((u_n)\) où, selon ce modèle, \(u_n\) est le nombre (en milliers) de voitures électriques immatriculées en France l'année \(2015+n\) avec \(n \in \mathbb{N}\). Ainsi, on a \(u_0 = 17,3\).
	
	\noindent
	\textbf{--- Second modèle :}
	
	On définit la suite \((v_n)\) par \(v_0 = 17,3\) et pour tout entier naturel \(n\), \(v_{n+1} = 0,7v_n +10\). D'après ce modèle et pour tout entier naturel \(n\), \(v_n\) est le nombre (en milliers) de voitures électriques immatriculées en France l'année \(2015+n\).
	
	\subsection*{a.}
	
	Ajouter 21 \%, c'est multiplier par \(1 + \dfrac{21}{100} = 1 + 0,21 = 1,21\). Donc :
	
	\begin{itemize}
		\item $u_1 = u_0 \times 1,21 = 17,3 \times 1,21 \approx 20,9$
		\item $u_2 \approx 25,3$
		\item $u_3 \approx 30,6$.
	\end{itemize}
	et :
	\begin{itemize}
		\item $v_1 = 0,7 \times 17,3 + 10 \approx 22,1$
		\item $v_2 \approx 25,5$
		\item $v_3 \approx 27,8$.
	\end{itemize}
	
	\subsection*{b.}
	
	Le second modèle donne des nombres plus proches de la réalité des ventes.
	
	\subsection*{2.}
	
	\subsubsection*{a.}
	
	On a vu que, quel que soit le naturel \(n\), \(u_{n+1} = 1,21u_n\) : ceci montre que la suite \((u_n)\) est une suite géométrique de raison 1,21 et de premier terme \(17,3\).
	
	\subsubsection*{b.}
	
	On sait que, quel que soit le naturel \(n\), \(u_n = 17,3 \times 1,21^n\).
	
	\subsubsection*{c.}
	
	On considère l'algorithme en langage Python ci-dessous.
	
	\begin{lstlisting}[style=python]
	u = 17.3
	n = 0
	while u < 50 :
		u = 1.21 * u
		n = n + 1
	\end{lstlisting}
	
	L'algorithme va donner \(n = 5\) pour s'arrêter avant que \(u\) ne dépasse 50, soit 50 000 voitures électriques.

\end{document}
